%=========================
% FrontMatter/summary_vi.tex
% Tóm tắt báo cáo bằng tiếng Việt
%=========================
\begin{center}
    \textbf{\Large TÓM TẮT}
\end{center}

Đề tài \textit{\enquote{Xây dựng Trợ lý Mua sắm Đa phương tiện sử dụng RAG và Computer Vision}} hướng đến việc phát triển một nền tảng mua sắm trực tuyến thông minh, ứng dụng trí tuệ nhân tạo nhằm nâng cao khả năng tìm kiếm, tư vấn và lựa chọn sản phẩm trong lĩnh vực thương mại điện tử. Hệ thống được xây dựng như một nền tảng tích hợp, nơi người dùng có thể tra cứu thông tin sản phẩm, thực hiện các nghiệp vụ mua sắm và tương tác với trợ lý bằng ngôn ngữ tự nhiên hoặc hình ảnh.

Trong quá trình triển khai, đề tài tập trung hiện thực các chức năng cốt lõi gồm tìm kiếm sản phẩm bằng văn bản, tìm kiếm bằng hình ảnh, tìm kiếm kết hợp văn bản--hình ảnh và tư vấn sản phẩm qua hội thoại nhiều lượt. Trợ lý được phát triển dựa trên kỹ thuật Retrieval-Augmented Generation (RAG), mô hình ngôn ngữ lớn và bộ nhớ hội thoại nhằm khai thác dữ liệu sản phẩm được truy xuất, duy trì ngữ cảnh và tạo phản hồi phù hợp với yêu cầu của người dùng. Bên cạnh đó, các mô hình biểu diễn đa phương thức và kỹ thuật thị giác máy tính được sử dụng để trích xuất đặc trưng, hỗ trợ truy hồi những sản phẩm tương đồng với hình ảnh đầu vào.

Trong phạm vi báo cáo, tác giả trình bày quá trình phân tích, thiết kế, xây dựng và tích hợp các thành phần của hệ thống, bao gồm Frontend, Backend nghiệp vụ, cơ sở dữ liệu MySQL, AI Service và cơ sở dữ liệu vector Qdrant; đồng thời mô tả cơ chế đồng bộ và phối hợp giữa các thành phần nhằm bảo đảm tính nhất quán của dữ liệu.

Tôi cũng tiến hành đánh giá các cấu hình truy hồi, vai trò của RAG và bộ nhớ hội thoại, cùng với việc kiểm thử chức năng, hiệu năng, bảo mật và khả dụng của hệ thống. Kết quả cho thấy giải pháp có khả năng hỗ trợ người dùng tìm kiếm sản phẩm bằng nhiều phương thức, cung cấp phản hồi có căn cứ và duy trì ngữ cảnh hội thoại trong phạm vi thử nghiệm. Hệ thống thể hiện khả năng ứng dụng thực tiễn của trí tuệ nhân tạo trong việc hỗ trợ quá trình tìm kiếm và ra quyết định mua sắm, đồng thời tạo cơ sở để tiếp tục cải thiện chất lượng truy hồi, khả năng xử lý ràng buộc và trải nghiệm người dùng.
